% Generated draft: requires mathematical and source review.
\documentclass[11pt]{article}
\usepackage[margin=25mm]{geometry}
\usepackage{fontspec}
\setmainfont{Latin Modern Roman}
\usepackage{amsmath,amssymb,mathtools,unicode-math}
\setmathfont{Latin Modern Math}
\usepackage{enumitem,needspace,xurl,hyperref}
\setlength{\parindent}{0pt}
\setlength{\parskip}{5pt}
\setlength{\emergencystretch}{3em}
\newcommand{\N}{\mathbb{N}}
\newcommand{\Z}{\mathbb{Z}}
\newcommand{\Q}{\mathbb{Q}}
\newcommand{\R}{\mathbb{R}}
\newcommand{\C}{\mathbb{C}}
\newcommand{\F}{\mathbb{F}}
\newcommand{\A}{\mathbb{A}}
\newcommand{\PP}{\mathbb{P}}
\newcommand{\E}{\mathbb{E}}
\newcommand{\eps}{\varepsilon}
\newcommand{\dd}{\,\mathrm{d}}
\begin{document}
\section*{UnsolvedMath ID: 800241}
\textbf{Source problem code:} catalogue entry
\subsection{Definitions and mathematical statement}
For interpreting the terminology, let \(n\geq 1\) and let \(\Omega\subset\C^n\) be a bounded domain, meaning a bounded, connected open set. The source does not specify the complex dimension. Call \(\Omega\) pseudoconvex if it admits a continuous plurisubharmonic exhaustion function \(\rho\colon\Omega\to\R\), so that \(\{z\in\Omega:\rho(z)<c\}\) is relatively compact in \(\Omega\) for every \(c\in\R\).

The symbol \(B\) is not defined in the supplied record. A possible interpretation is the Bergman metric. Under this interpretation, define
\[
A^2(\Omega)=\left\{f\colon\Omega\to\C:\ f\text{ is holomorphic and }\int_\Omega |f|^2\,\dd V<\infty\right\},
\qquad
K_\Omega(z,z)=\sup_{\substack{f\in A^2(\Omega)\\ \|f\|_{L^2(\Omega)}\leq 1}}|f(z)|^2.
\]
Its metric tensor is
\[
g_{j\bar k}(z)=\frac{\partial^2\log K_\Omega(z,z)}{\partial z_j\,\partial\bar z_k},
\qquad
B_\Omega(z;v)^2=\sum_{j,k=1}^n g_{j\bar k}(z)v_j\overline{v_k}.
\]
Let \(d_{B,\Omega}\) be the distance obtained by taking the infimum of the \(B_\Omega\)-lengths of piecewise continuously differentiable curves joining two points of \(\Omega\). Completeness means that every \(d_{B,\Omega}\)-Cauchy sequence converges in this distance to a point of \(\Omega\).

\textbf{Source question, with unresolved qualifications.} Is \(B\) complete on bounded pseudoconvex domains with boundary of regularity ``less than \(C^1\)''? If \(B\) means the Bergman metric and \(\mathcal D\) denotes the intended class of such domains, the proposed assertion is
\[
\text{for every }\Omega\in\mathcal D,\quad (\Omega,d_{B,\Omega})\text{ is complete}.
\]
The supplied record neither identifies \(B\) nor defines \(\mathcal D\): it gives no precise boundary regularity threshold and does not specify whether all boundary regularity assumptions are to be dropped. Thus this conditional formulation requires clarification before it becomes a definite mathematical assertion.
\subsection{Short English statement}
The question asks whether bounded pseudoconvex domains with boundary regularity below \(C^1\) have a complete metric denoted by \(B\). If \(B\) is the Bergman metric, completeness means that every Cauchy sequence for its induced distance converges inside the domain. The metric notation, intended boundary class, and dimension require clarification.
\subsection{Sources}
\begin{itemize}
\item Local UnsolvedMath registry: permanent ID \(800241\), exact source code \texttt{catalogue entry}, title \emph{Several Complex Variables Problem List --- Problem 38}. The supplied registry labels the record open; this label is not an independently verified mathematical status.
\item Buzzard, Fornaess, and Krantz, \emph{The Several Complex Variables Problem List} (1995), Problem 38. Supplied source URL: \url{https://arxiv.org/abs/math/9509202}. The registry attributes the collection to these authors; the question itself is signed J. McNeal, with supplied contact \texttt{mcneal@math.princeton.edu}.
\item Supplied provenance records access on 2026-10-16 and extraction by \texttt{arxiv\_latex}. These are supplied provenance statements, not a claim of independent source access or verification for this draft.
\end{itemize}
\end{document}
